
\chapter{Fundamentação Teórica}

\section{Hardware: Unitree Go2}
O Unitree Go2 integra um hardware robusto com tecnologia de ponta em atuadores e 
sensoriamento. A versão EDU, em particular, destaca-se por oferecer ferramentas 
avançadas para adaptação de software e integração com sistemas externos, 
expandindo as capacidades nativas do robô \cite{unitree2024go2}.

Como um robô quadrúpede, o Go2 possui uma estrutura inspirada na biomecânica 
animal, com um tronco central e quatro membros simétricos. O robô tem 12 graus de 
liberdade (DoF), distribuídos em três juntas por perna:

\begin{itemize}
    \item Quadril (Hip): Responsável pelo movimento de abdução e adução (abrir 
    e fechar a perna lateralmente).
    \item Coxa (Thigh): Realiza a protração e retração (movimento de avanço e 
    recuo).
    \item Panturrilha (Calf): Executa a flexão e extensão, equivalente ao 
    movimento do joelho.
\end{itemize}

Para interagir com o ambiente de forma autônoma e segura, o Go2 utiliza um 
conjunto sensorial de alta tecnologia:

\begin{itemize}
    \item Lidar 4D L1: Este sensor de medição a laser de alta velocidade realiza 
    até 21.600 amostras por segundo. Ele fornece informações do ambiente em tempo 
    real, permitindo o reconhecimento preciso de terrenos e estruturas.
    \item Câmera Grande Angular HD: Capaz de transmitir vídeo em  até 1080p a 
    30fps, oferece uma visão rica em detalhes e baixa latência para monitoramento 
    ou processamento de visão computacional.
    \item Sensores de Força nas Patas: Localizados nas extremidades de cada 
    membro, esses sensores detectam o contato com o solo e a intensidade da 
    pressão exercida. Esses dados são cruciais para o ajuste de equilíbrio em 
    superfícies irregulares.
    \item IMU de 9 Eixos: A Unidade de Medida 
    Inercial combina três sensores fundamentais:
    \begin{itemize}
        \item Acelerômetro: Mede a aceleração linear nos eixos X, Y e Z.
        \item Giroscópio: Monitora a velocidade angular (rotação) sobre os eixos.
        \item Magnetômetro: Identifica a orientação em relação ao campo magnético terrestre.
    \end{itemize}
\end{itemize}

\section{Estrutura do estimador de estados MUSE}
O MUSE (Multi-Sensor State Estimator) representa uma evolução significativa em relação ao 
seu antecessor apresentado no trabalho "Proprioceptive Sensor Fusion for Quadruped Robot 
State Estimation" \cite{fink2020proprioceptive}.Como um estimador de estados de última geração, 
ele introduz a fusão de dados exteroceptivos (informações do ambiente) com 
sensores proprioceptivos, além de incorporar um módulo dedicado à detecção de 
deslize em um ou mais membros do robô. Essa arquitetura integrada permite que o 
MUSE apresente um drift (desvio) de posição reduzido, mantendo uma robustez 
maior, mesmo em cenários de falha sensorial ou navegação em terrenos altamente 
irregulares e imprevisíveis \cite{nistico2025muse}.

\subsection{Odometria visual e por LiDAR}
Para o processamento das odometrias de câmera e do Lidar, o MUSE utiliza o 
algoritmo KISS-ICP. A escolha desta ferramenta justifica-se por sua versatilidade 
e eficiência computacional. O KISS-ICP demonstra uma capacidade ótima  de 
adaptação a diversos contextos operacionais utilizando um conjunto reduzido de 
parâmetros, o que facilita o ajuste do sistema para diferentes ambientes sem a 
necessidade de muitas mudanças no setup das variáveis \cite{vizzo2023kiss}.

\subsection{Odometria das pernas (leg odometry)}
A odometria das pernas estima o deslocamento do tronco em tempo real através da 
cinemática direta dos membros. Para realizar esse processo, parte-se da premissa 
que as patas em fase de apoio (contato com o solo) estão estáticas em relação ao 
mundo, permitindo calcular a movimentação do corpo a partir dos ângulos das 
juntas. No entanto, para que essa estimativa seja correta e precisa, assume-se 
inicialmente a ausência de deslizamento.

\subsection{Módulo de detecção de deslize (slip detection)}
Em superfícies com baixo coeficiente de atrito, o robô está sujeito a deslizamentos 
que geram erros acumulativos (drift) na posição estimada. Para robôs com pernas, 
uma metodologia de detecção baseada na cinemática é mais adequada do que metodologias 
baseadas em força/torque, uma vez que a detecção de deslizamentos por forças exigiria
 o uso de sensores de força/torque de 6 eixos \cite{nistico2022slip}.
O MUSE soluciona esse 
problema através de um módulo de detecção de deslize que monitora dois parâmetros: 
a discrepância entre a velocidade/posição desejada e a atingida. Quando o desvio 
entre esses valores ultrapassa um limiar crítico, o sistema identifica o deslize 
em uma pata específica. Essa detecção é vital não apenas para a precisão da 
estimativa, mas também para permitir que a estratégia de controle do robô se 
adapte instantaneamente, buscando evitar erros na posição estimada. \cite{nistico2022slip}

\subsection{Observador de atitude}
A orientação do robô é processada por uma estrutura em cascata, 
composta por um Observador Não Linear (NLO) em série com um Filtro de Kalman 
Exógeno (XKF) \cite{fink2020proprioceptive}.

\begin{itemize}
    \item Observador Não Linear (NLO): Atua na observação inicial do sistema, 
    fornecendo uma estimativa rápida da atitude. Embora eficiente, o NLO pode 
    divergir se a estimativa inicial estiver muito distante do estado real.
    \item Filtro de Kalman Exógeno (XKF): Entra como uma camada de correção, 
    garantindo a estabilidade global do sistema. Ele combina as propriedades de 
    convergência quase-ótimas do Filtro de Kalman com a robustez do NLO, 
    garantindo que a orientação seja mantida de forma precisa mesmo em condições 
    dinâmicas.
\end{itemize}

A escolha do XKF em vez do tradicional EKF (Filtro de Kalman Estendido) foi 
proposital. O XKF consegue se manter estável em situações onde o EKF apresenta 
instabilidade, garantindo que o sistema seja robusto o tempo todo, como exemplificado 
na obra "The eXogenous Kalman Filter (XKF)" \cite{johansen2017xkf}.


\subsection{Fusão de sensores}
O MUSE realiza a fusão das medidas inerciais (IMU) com a odometria das pernas e 
os dados exteroceptivos (câmera/Lidar) por meio de algoritmos baseados em filtros 
de Kalman. Um diferencial é o desacoplamento entre 
atitude e posição, o que transforma a dinâmica no Filtro de Kalman em um sistema 
linear e invariante no tempo (LTI), o que como já citado anteriormente, garantindo a convergência matemática do filtro.

Além disso, o algoritmo lida de forma inteligente com a latência do lidar ou 
câmera. Enquanto os sensores exteroceptivos (que possuem frequência de 
processamento mais baixa) não enviam novos dados, o MUSE utiliza a alta frequência
da IMU e da odometria das pernas para manter uma estimativa imediata. Assim que os 
dados da câmera ou Lidar chegam, o sistema aplica correções, 
ajustando a trajetória estimada e eliminando erros acumulados durante o intervalo
de processamento \cite{nistico2025muse}.

\section{validação em ambiente simulado}
A primeira parte de validação do estimador de estados será conduzida em ambiente 
simulado, utilizando o simulador Gazebo. Este software permite uma representação
muito próxima à realidade da cinemática e da dinâmica do robô Go2. Além de 
emular a física do robô, o Gazebo também possibilita a modelagem dos sensores. 
Para trazer maior confiabilidade à validação do MUSE, no gazebo será configurado 
ruído nos sensores, aproximando as leituras virtuais das imperfeições e 
incertezas presentes no mundo físico.

Para a avaliação precisa da performance do estimador e do comportamento do 
sistema, a simulação contará com ferramentas específicas de análise de dados:

\begin{itemize}
    \item PlotJuggler: O monitoramento e a análise dos dados brutos serão 
    realizados através do PlotJuggler. Trata-se de uma ferramenta avançada de 
    visualização de séries temporais, essencial para a leitura de variáveis do 
    sistema. Ele será empregado principalmente para analisar o comportamento 
    temporal das leituras da IMU e avaliar a resposta dos sinais de controle.
    \item Point Cloud Library (PCL): Para o processamento tridimensional e a 
    visualização das nuvens de pontos geradas pelo Lidar simulado, bem como para 
    a representação  dos mapas espaciais construídos, será utilizada a PCL 
    (Point Cloud Library). Esta biblioteca oferece uma ampla gama de algoritmos e 
    renderizadores que permitem a inspeção detalhada da geometria do ambiente e a 
    verificação visual da precisão da odometria exteroceptiva do robô.
\end{itemize}

